\chapter{Sistema de Negocio}
\section{Descripción del Sistema de Negocio}
El sistema de negocio del proceso de “Concurso de Personal” en la empresa Tambo comprende el conjunto de actividades, actores y recursos involucrados en la gestión del reclutamiento y selección de personal. Este sistema inicia con la identificación de la necesidad de contratación y la publicación de vacantes, continúa con la recepción y evaluación de postulantes, y culmina con la selección y contratación del candidato idóneo.

En este contexto, participan diversos actores como el postulante, el reclutador y el área de recursos humanos, quienes interactúan mediante procesos definidos para garantizar la incorporación eficiente de talento humano a la organización.

El sistema actual presenta limitaciones relacionadas con la alta carga operativa, procesos manuales y falta de integración de la información, lo que genera retrasos y posibles pérdidas de talento.

\section{Objetivos del Sistema}
\textbf{Objetivo General}
Diseñar e implementar un sistema de información que optimice el proceso de gestión del Concurso de Personal 
en la empresa Tambo, reduciendo el Tiempo de contratación de 15-20 días a un máximo de 7 días hábiles 
y la tasa de abandono de candidatos aptos del $35\%$ al $10\%$.

\textbf{Objetivos Específicos}
\begin{itemize}
    \item Implementar un módulo de filtrado automático que pre-califique el $100\%$ de las postulaciones 
    en menos de 5 segundos por candidato, sin intervención manual del Analista.
    \item Consolidar la información de postulantes de todas las tiendas en una única base de datos, 
    eliminando el uso de hojas de cálculo locales en un plazo de 3 meses desde la implementación.
    \item Alcanzar un Tiempo de contratación promedio de 5-7 días hábiles para perfiles operativos 
    operativos en los primeros 6 meses de operación del sistema.
    \item Permitir al Administrador de Tienda visualizar el estado actualizado de sus candidatos en menos de 5 
    segundos desde cualquier dispositivo, sin necesidad de contactar al Analista.
    \item Facilitar la toma de decisiones mediante información organizada y accesible.
\end{itemize}

\section{Alcance del Sistema}

El sistema abarcará las siguientes funcionalidades:

\begin{itemize}
    \item Publicación de convocatorias laborales en el portal del sistema.
    \item Registro y almacenamiento de postulaciones con validación de datos.
    \item Pre-calificación automática de candidatos mediante filtros lógicos configurables.
    \item Gestión de estados del candidato a lo largo de todo el pipeline.
    \item Selección y contratación de personal.
\end{itemize}

El sistema no incluirá:
\begin{itemize}
    \item Procesos de inducción, onboarding o capacitación interna del personal contratado.
    \item Gestión de desempeño posterior a la contratación (evaluaciones, KPIs del colaborador).
    \item Procesos de nómina, liquidaciones o gestión salarial.
\end{itemize}